\section{High Pass Filter}\label{sec:hpf}
Ngược với bộ lọc thông thấp, bộ lọc thông cao chỉ cho qua tín hiệu có tần số lớn hơn tần số cắt dưới $f_L$. Tài liệu tham khảo:
\begin{center}\footnotesize\url{https://www.allaboutcircuits.com/video-tutorials/op-amps-low-pass-and-high-pass-active-filters/}\end{center}

Yêu cầu: tự thiết kế mạch high pass filter trên PSpice (đề xuất cấu hình đảo hoặc không đảo), mô phỏng và giải thích hoạt động.

\subsection{Thiết kế}
\begin{notebox}{Cấu hình đã chọn và công thức tần số cắt}
\dotlines{4}
\end{notebox}
\calc[6]
\begin{table}[H]
\centering
\begin{tabular}{l c}
\toprule
\textbf{Thông số} & \textbf{Giá trị}\\
\midrule
$R_1$, $R_2$ & \fillcell\\
$C$ & \fillcell\\
$f_L$ (lý thuyết) & \fillcell\\
Gain băng thông (V/V) & \fillcell\\
\bottomrule
\end{tabular}
\caption{Thông số high pass filter}
\end{table}

\subsection{Mô phỏng}
\simcfg{AC Sweep}{ghi rõ kiểu sweep, tần số đầu/cuối, số điểm}
\simresult
\simfig{ex06_sch}{Sơ đồ mạch high pass filter (PSpice)}
\simfig{ex06_ac}{Kết quả AC Sweep của high pass filter}

\begin{table}[H]
\centering
\begin{tabular}{l c c}
\toprule
\textbf{Đại lượng} & \textbf{Lý thuyết} & \textbf{Mô phỏng}\\
\midrule
$f_L$ (Hz) & \fillcell & \fillcell\\
Gain băng thông (dB) & \fillcell & \fillcell\\
\bottomrule
\end{tabular}
\caption{So sánh lý thuyết -- mô phỏng}
\end{table}

\begin{notebox}{Giải thích hoạt động của mạch}
\dotlines{5}
\end{notebox}
